\documentclass[12pt,a4paper]{article}

\usepackage[russian]{babel}
\usepackage{geometry}
\usepackage{booktabs}
\usepackage{longtable}
\usepackage{enumitem}
\usepackage{hyperref}
\usepackage{titlesec}
\usepackage{float}
\usepackage{amsmath}
\usepackage{indentfirst}

\geometry{left=2.5cm, right=2.5cm, top=2cm, bottom=2cm}

\titleformat{\section}{\Large\bfseries\filcenter}{\thesection}{1em}{}
\titleformat{\subsection}{\large\bfseries\filcenter}{\thesubsection}{1em}{}

\setlength{\parindent}{1.25cm}
\setlength{\parskip}{0.5em}

\setlist{noitemsep, topsep=0pt}

\title{
  \textbf{ЭКОНОМИЧЕСКАЯ СТРАТЕГИЯ} \\
  \vspace{0.3cm}
  \textbf{Развития образовательного бизнеса} \\
  \vspace{0.3cm}
  \textbf{на 12 месяцев} \\
  \vspace{0.3cm}
  \large
  для двух партнёров---самозанятых (НПД) \\
  \vspace{0.5cm}
}

\author{Аналитический отдел}
\date{Сентябрь 2026 г.}

\begin{document}

\maketitle

\newpage
\tableofcontents
\newpage

% ========================================================================
% ВВЕДЕНИЕ
% ========================================================================
\section*{Введение}
\addcontentsline{toc}{section}{Введение}

Данный документ представляет собой детальную экономическую стратегию развития образовательного бизнеса на 12 месяцев (октябрь 2026 --- сентябрь 2027 гг.) для двух партнёров---самозанятых в возрасте 15---16 лет. Бизнес основывается на преподавании математики, физики и машинного обучения (ML) для учащихся 1---9 классов по формату групповых занятий (4---8 человек) в онлайн- и офлайн-формате.

\textbf{Цель стратегии:} построить устойчивый образовательный бизнес с доходом 600\,000---800\,000 руб./партнёр/год к концу первого года, с перспективой масштабирования до уровня образовательной компании (ориентир --- МШП) за 3---5 лет.

\textbf{Ключевые параметры стартапа:}

\begin{itemize}
  \item Стартовый капитал: 10\,000 руб.
  \item Статус: самозанятые (НПД, ставка 4\%).
  \item Возраст партнёров: 15---16 лет.
  \item Ограничение дохода по НПД: 2,4 млн руб./год на одного партнёра.
  \item Невозможность трудоустраивать сотрудников по трудовому договору.
  \item Бесплатное помещение в лицее для офлайн-заняний.
  \item Цена: 750---1000 руб./час/человек.
\end{itemize}

Все расчёты и прогнозы основаны на реалиях российского рынка частного образования, действующем законодательстве о самозанятости (ФЗ №422-ФЗ) и типовых доходных нагрузках для подросткового возраста.

% ========================================================================
% АНАЛИЗ РЫНКА
% ========================================================================
\section{Анализ рынка}

\subsection{Общие размеры и динамика}

Российский рынок частного образования пережил значительные изменения после 2022 года. Международные EdTech-платформы (Coursera, Udemy, edX) сократили операции в РФ, что создало нишу для локальных игроков.

\begin{table}[H]
\centering
\caption{Ключевые параметры российского рынка частного образования}
\label{tab:market_overview}
\begin{tabular}{ll}
\toprule
Параметр & Значение \\
\midrule
Объём рынка частного образования (2026) & 150---200 млрд руб. \\
Годовой рост (2023---2026, средний) & 8---12\% \\
Доля частных репетиторов & ${>}$ 60\% \\
Доля онлайн---платформ & ${<}$ 25\% \\
Доля очных школ/клубов & ${>}$ 70\% \\
Средняя цена 1 часа репетиторства & 700---1\,500 руб. \\
\bottomrule
\end{tabular}
\end{table}

\textbf{Вывод:} Рынок остаётся фрагментированным. Крупные игроки (МШП, Яндекс.Практикум, Skillbox) занимают премиум-сегмент, но спрос на доступное и качественное образование остаётся высоким.

\subsection{Конкурентный анализ}

\subsubsection{Московская Школа Программирования (МШП)}

МШП --- лидер в премиум-образовании по программированию и ML.

\begin{table}[H]
\centering
\caption{Конкурентный анализ --- МШП}
\label{tab:msp_analysis}
\begin{tabular}{ll}
\toprule
Параметр & МШП \\
\midrule
Год основания & 2015 \\
Количество преподавателей & ${\approx}$ 800 \\
Средняя цена & 3\,000---5\,000 руб./час \\
Формат & Онлайн + офлайн \\
ЦА & Дети 7---17 лет, взрослые \\
Сильные стороны & Бренд, курсы ML, олимпиады \\
Слабые стороны & Высокая цена, отток клиентов \\
\bottomrule
\end{tabular}
\end{table}

\subsubsection{Код будущего (Сбер)}

Государственная инициатива по цифровому образованию.

\begin{table}[H]
\centering
\caption{Конкурентный анализ --- Код будущего}
\label{tab:code_future}
\begin{tabular}{ll}
\toprule
Параметр & Код будущего \\
\midrule
Формат & Онлайн-курсы \\
Цена & Бесплатно / платно (сертификат) \\
ЦА & Школьники 10---11 классов \\
Сильные стороны & Господдержка, масштаб \\
Слабые стороны & Общий характер, нет персонального подхода \\
\bottomrule
\end{tabular}
\end{table}

\subsubsection{Крупные EdTech-платформы}

\begin{table}[H]
\centering
\caption{Конкурентный анализ --- крупные платформы}
\label{tab:edtech_competitors}
\begin{tabular}{llll}
\toprule
Платформа & Формат & Цена & Особенности \\
\midrule
Яндекс.Практикум & Онлайн & 15\,000---30\,000/мес & Взрослая аудитория \\
Skillbox & Онлайн & 12\,000---25\,000/мес & Профессиональная переподготовка \\
GeekBrains & Онлайн+офлайн & 10\,000---20\,000/мес & Широкий спектр курсов \\
Нетология --- Онлайн & 15\,000---25\,000/мес & Бизнес-образование \\
\bottomrule
\end{tabular}
\end{table}

\subsubsection{Частные репетиторы}

Частные репетиторы доминируют на рынке (${>}$ 60\% доли), но конкурируют между собой по цене и качеству. Средняя цена: 700---1\,500 руб./час.

\subsection{Выводы по конкурентному анализу}

\begin{enumerate}
  \item \textbf{Ценовой разррыв:} Премиум-игроки (МШП, Яндекс.Практикум) предлагают курсы за 3\,000+ руб./час, создавая спрос на доступные ($<$ 1\,000 руб.) но качественные аналоги.
  \item \textbf{ML как ниша:} В 2026 г. менее 5\% школьных репетиторов предлагают ML. Это ключевое конкурентное преимущество.
  \item \textbf{Возрастной фактор:} Подростковые репетиторы (${approx}$ 16 лет) ближе к ученикам 1---9 классов, чем взрослые преподаватели.
  \item \textbf{Фрагментация рынка:} Мало профессионально организованных малых игроков, занимающихся ML для школьников.
\end{enumerate}

% ========================================================================
% ЭКОНОМИЧЕСКАЯ МОДЕЛЬ
% ========================================================================
\section{Экономическая модель}

\subsection{Модель доходов}

Доход бизнеса складывается из платы за обучение по трём направлениям: ML, математика, физика.

\begin{equation}
R = \sum_{i=1}^{n} S_i \cdot H_i \cdot P_i \cdot W
\label{eq:revenue}
\end{equation}

Где:
\begin{itemize}
  \item $R$ --- месячный доход, руб.;
  \item $S_i$ --- количество учеников в группе $i$;
  \item $H_i$ --- часы занятий в неделю в группе $i$;
  \item $P_i$ --- цена за час на одного ученика, руб.;
  \item $W$ --- количество недель в месяце (${approx}$ 4.3).
\end{itemize}

\textbf{Пример расчёта:} 1 группа ML, 6 учеников, 2 часа/неделя, 850 руб./час:

\begin{equation}
R = 6 \cdot 2 \cdot 850 \cdot 4.3 = 43\,860 \text{ руб./мес}
\end{equation}

\subsection{Модель расходов}

\begin{table}[H]
\centering
\caption{Структура расходов (ежемесячно)}
\label{tab:expenses}
\begin{tabular}{llr}
\toprule
Статья расходов & Описание & Сумма, руб. \\
\midrule
Интернет + связь & Широкополосный доступ, мобильная связь & 1\,000 \\
Маркетинг & Таргет, печатные материалы & 2\,000 \\
Материалы & Блокноты, маркеры, расходники & 500 \\
Инструменты & ПО, онлайн-сервисы (до старта) & 1\,000 \\
Прочее & Мелкие непредвиденные расходы & 500 \\
\midrule
\textbf{Итого} & & \textbf{5\,000} \\
\bottomrule
\end{tabular}
\end{table}

\subsection{Налоговая модель (НПД)}

Самозанятые платят налог на процессуальный доход (НПД) по ставке 4\%. Применяется стандартный налоговый вычет в размере 400\,000 руб./год.

\begin{equation}
\text{Налог} = 0.04 \cdot \max(0,\ (D_{\text{год}} - 400\,000))
\label{eq:npd_tax}
\end{equation}

Где $D_{\text{год}}$ --- годовой доход одного партнёра.

\begin{table}[H]
\centering
\caption{Расчёт налога по НПД для разных уровней дохода}
\label{tab:npd_calculation}
\begin{tabular}{lrr}
\toprule
Годовой доход, руб. & Налог, руб. & Чистый доход, руб. \\
\midrule
500\,000 & 0 & 500\,000 \\
700\,000 & 12\,000 & 688\,000 \\
1\,000\,000 & 24\,000 & 976\,000 \\
1\,500\,000 & 44\,000 & 1\,456\,000 \\
2\,400\,000 & 80\,000 & 2\,320\,000 \\
\bottomrule
\end{tabular}
\end{table}

\textbf{Важно:} Доход свыше 2,4 млн руб./год облагается обычным подоходным налогом (13\%), что делает бизнес с НПД невыгодной для масштабирования. Максимальный доход по НПД: ${\approx}$ 2,32 млн руб./год на одного партнёра.

\subsection{Модель чистой прибыли}

\begin{equation}
\Pi = R_{\text{год}} \cdot \frac{1}{2} - E_{\text{год}} - T
\label{eq:net_profit}
\end{equation}

Где:
\begin{itemize}
  \item $\Pi$ --- чистая прибыль одного партнёра, руб.;
  \item $R_{\text{год}}$ --- годовой доход бизнеса, руб.;
  \item $E_{\text{год}}$ --- годовые расходы, руб. (${approx}$ 60\,000);
  \item $T$ --- налог по НПД, руб.
\end{itemize}

\subsection{Анализ точки безубыточности}

\begin{equation}
R_{\text{мес}} = E_{\text{мес}} / (1 - \alpha)
\end{equation}

Где $\alpha$ --- доля налога (4\% при доходе ${>}$ 400\,000/год).

При месячных расходах 5\,000 руб., точка безубыточности достигается при доходе ${\approx}$ 5\,200 руб./мес на одного партнёра (${approx}$ 10\,400 руб. в месяц — 2 группы с 3---4 учениками каждая).

% ========================================================================
% ЭТАПЫ РАЗВИТИЯ НА 12 МЕСЯЦЕВ
% ========================================================================
\section{Этапы развития на 12 месяцев}

\subsection{Этап 1: Подготовка (октябрь 2026)}

\subsubsection*{Цель:}
Завершить ВсОШ, зарегистрироваться как самозанятые, разработать учебные программы, создать рекламные материалы.

\subsubsection*{Действия:}
\begin{enumerate}[label=\arabic*.]
  \item Завершить подготовку к ВсОШ (до 12 октября).
  \item Зарегистрироваться как самозанятые через мобильное приложение ``Мой налог'' (требуется согласие родителей/законных представителей --- форма СЗ --- 21/2023).
  \item Разработать учебные программы:
    \begin{itemize}
      \item ML для 1---9 классов: 10 уроков, каждый --- 2 часа, включает Python, data science, простые нейросети.
      \item Математика для 1---9 классов: программа по ФОПФУ, олимпиадные задачи.
      \item Физика для 5---9 классов: основы механики, электроники, оптики.
    \end{itemize}
  \item Создать рекламные материалы (презентация, видео---превью, посты для соцсетей).
  \item Открыть Telegram-канал и Instagram-аккаунт.
  \item Подготовить договор с родителями учеников (образец --- приложение).
\end{enumerate}

\subsubsection*{Бюджет:}
\begin{table}[H]
\centering
\caption{Этап 1: Бюджет ``Подготовка''}
\begin{tabular}{lr}
\toprule
Статья & Сумма, руб. \\
\midrule
Материалы (бумага, канцелярия) & 2\,000 \\
Регистрация, справки & 500 \\
Маркетинг (создание контента) & 3\,000 \\
Прочее & 2\,500 \\
\midrule
\textbf{Итого} & \textbf{8\,000} \\
\bottomrule
\end{tabular}
\end{table}

\subsubsection*{Ожидаемый доход:} 0 руб.

\subsubsection*{Риски:}
\begin{itemize}
  \item \textbf{Задержка регистрации:} Самозанятость для несовершеннолетних требует дополнительного подтверждения --- риск 1---2 недели. \textbf{Митигация:} Подать документы за 2 недели до планового старта.
  \item \textbf{Сложности с курендурой:} ML --- специфичный продукт. \textbf{Митигация:} Сначала использовать проверенные open-source проекты (Kaggle, Google Colab).
\end{itemize}

\subsection{Этап 2: Первая группа (ноябрь --- декабрь 2026)}

\subsubsection*{Цель:}
Запустить первую пилотную группу ML и первую группу математики, проверить модель продаж и обратную связь.

\subsubsection*{Действия:}
\begin{enumerate}
  \item Запустить пилотную группу ML (4---6 учеников, 2 часа/неделя, 800 руб./час).
  \item Запустить первую группу математики (5 учеников, 2 часа/неделя, 850 руб./час).
  \item Собирать обратную связь и дорабатировать программы.
  \item Запустить реферальную программу (500 руб. за приглашение, 500 руб. на счёт ученика).
\end{enumerate}

\subsubsection*{Бюджет:}
\begin{table}[H]
\centering
\caption{Этап 2: Бюджет ``Первые группы''}
\begin{tabular}{lr}
\toprule
Статья & Сумма, руб. \\
\midrule
Маркетинг (таргетинг в VK/Telegram) & 3\,000 \\
Материалы (ноутбуки, канцелярия) & 1\,500 \\
Прочее & 1\,000 \\
\midrule
\textbf{Итого} & \textbf{5\,500} \\
\bottomrule
\end{tabular}
\end{table}

\subsubsection*{Ожидаемый доход:}
\begin{table}[H]
\centering
\caption{Этап 2: Доход ``Первые группы''}
\begin{tabular}{lr}
\toprule
Группа & Доход, руб. \\
\midrule
ML (5 человек × 2 ч × 800 × 4.3 нед) & 34\,400 \\
Математика (5 человек × 2 ч × 850 × 4.3 нед) & 36\,550 \\
\midrule
\textbf{Итого} & \textbf{70\,950} \\
\bottomrule
\end{tabular}
\end{table}

Чистая прибыль на партнёра: ${\approx}$ 28\,000 руб./мес.

\subsubsection*{Риски:}
\begin{itemize}
  \item \textbf{Недостаток спроса:} В начале может быть сложно набрать группу. \textbf{Митигация:} Использовать собственную сеть знакомых, родственников, предложить пробное занятие бесплатно.
  \item \textbf{Сложности с платёжкой:} Родители могут не доверять подростковому репетитору. \textbf{Митигация:} Предлагать пробное занятие, отзывы реальных учеников, работать через проверенную площадку (если есть возможность).
\end{itemize}

\subsection{Этап 3: Масштабирование (январь --- март 2027)}

\subsubsection*{Цель:}
Увеличить количество групп до 4---5, добиться стабильного спроса, оптимизировать расписание.

\subsubsection*{Действия:}
\begin{enumerate}
  \item Добавить вторую группу ML (цель: 5---6 учеников).
  \item Добавить вторую группу математики.
  \item Добавить группу физики (4---5 учеников).
  \item Оптимизировать расписание: чередовать онлайн/офлайн.
  \item Автоматизировать запись через Telegram-бота (бесплатно через @ManyChatBot или аналог).
\end{enumerate}

\subsubsection*{Бюджет:}
\begin{table}[H]
\centering
\caption{Этап 3: Бюджет ``Масштабирование''}
\begin{tabular}{lr}
\toprule
Статья & Сумма, руб./мес \\
\midrule
Маркетинг & 3\,000 \\
Материалы & 1\,000 \\
Инструменты (бот, ПО) & 1\,000 \\
Прочее & 1\,000 \\
\midrule
\textbf{Итого} & \textbf{6\,000} \\
\bottomrule
\end{tabular}
\end{table}

\subsubsection*{Ожидаемый доход:}
\begin{table}[H]
\centering
\caption{Этап 3: Доход ``Масштабирование'' (март)}
\begin{tabular}{ll}
\toprule
Группа & Доход, руб./мес \\
\midrule
ML группа 1 (6 человек × 2 ч × 850) & 43\,860 × 4.3 = 188\,598 \\
ML группа 2 (5 человек × 2 ч × 850) & 36\,550 × 4.3 = 157\,165 \\
Математика 1 (6 человек × 2 ч × 900) & 43\,200 × 4.3 = 185\,760 \\
Математика 2 (5 человек × 2 ч × 900) & 36\,000 × 4.3 = 154\,800 \\
Физика (4 человека × 2 ч × 850) & 27\,200 × 4.3 = 116\,960 \\
\midrule
\textbf{Итого} & \textbf{803\,283} \\
\bottomrule
\end{tabular}
\end{table}

Доход на партнёра: ${\approx}$ 400\,000 руб./мес (пик).

\subsubsection*{Риски:}
\begin{itemize}
  \item \textbf{Перегрузка:} 4 группы × 2 часа = 8 часов/неделя на одного партнёра. \textbf{Митигация:} Распределять нагрузку между партнёрами, использовать ассистентов (по договору ГПХ --- если позволяет возраст).
  \item \textbf{Потеря учеников:} В начале учебного года возможен отток. \textbf{Митигация:} Предлагать пробные занятия, скидки на первый месяц.
\end{itemize}

\subsection{Этап 4: Весенний период и стабилизация (апрель --- июнь 2027)}

\subsubsection*{Цель:}
Стабилизировать доход, подготовить к экзаменам, сохранить здоровье участников.

\subsubsection*{Действия:}
\begin{enumerate}
  \item Сократить интенсивность занятий (1.5---2 часа/неделя на группу вместо 2).
  \item Предложить экзаменационные зачёты и подготовительные курсы.
  \item Провести анкетирование учеников для улучшения программ.
  \item Начать работу над летними лагерями/ML-курсами.
\end{enumerate}

\subsubsection*{Бюджет:}
\begin{table}[H]
\centering
\caption{Этап 4: Бюджет ``Весенний период''}
\begin{tabular}{lr}
\toprule
Статья & Сумма, руб./мес \\
\midrule
Маркетинг & 2\,000 \\
Материалы & 1\,000 \\
Прочее & 1\,000 \\
\midrule
\textbf{Итого} & \textbf{4\,000} \\
\bottomrule
\end{tabular}
\end{table}

\subsubsection*{Ожидаемый доход:} ${\approx}$ 150\,000---200\,000 руб./мес (сокращение из---за экзаменов).

\subsubsection*{Риски:}
\begin{itemize}
  \item \textbf{Снижение дохода в экзаменационный период.} \textbf{Митигация:} Фокус на подготовку к экзаменам и олимпиадам как платный сервис.
  \item \textbf{Выгорание.} \textbf{Митигация:} Планировать перерывы, чередовать группы.
\end{itemize}

\subsection{Этап 5: Лето (июль --- август 2027)}

\subsubsection*{Цель:}
Предложить интенсивные летние курсы ML, поддержать доход в низкий сезон.

\subsubsection*{Действия:}
\begin{enumerate}
  \item Запустить летний интенсив ML (10 учеников, 3 часа/неделя, 900 руб./час).
  \item Предложить гибрый график (онлайн/офлайн).
  \item Собрать кейсы и портфолио учеников.
\end{enumerate}

\subsubsection*{Ожидаемый доход:}
\begin{table}[H]
\centering
\caption{Этап 5: Доход ``Лето''}
\begin{tabular}{lr}
\toprule
Программа & Доход, руб./мес \\
\midrule
Летний ML интенсив (10 человек × 3 ч × 900 × 4.3) & 116\,100 \\
Регулярные группы (сокращённые) & 80\,000 \\
\midrule
\textbf{Итого} & \textbf{196\,100} \\
\bottomrule
\end{tabular}
\end{table}

\subsection{Этап 6: Новый учебный год (сентябрь 2027)}

\subsubsection*{Цель:}
Вернуться к полной загрузке, запустить новые группы, подвести итоги первого года.

\subsubsection*{Действия:}
\begin{enumerate}
  \item Запустить 4---5 групп с новыми учениками.
  \item Представить обновлённые программы.
  \item Подвести итоги года (отчёт, планы на следующий год).
\end{enumerate}

\subsubsection*{Ожидаемый доход:} ${\approx}$ 200\,000---250\,000 руб./мес на партнёра.

% ========================================================================
% ЮРИДИЧЕСКАЯ СХЕМА
% ========================================================================
\section{Юридическая схема}

\subsection{Регистрация самозанятых несовершеннолетних}

В России несовершеннолетние (до 18 лет) могут быть зарегистрированы как самозанятые с 2021 года (ФЗ №422-ФЗ, ст. 332.17). Порядок:

\begin{enumerate}
  \item Законный представитель (родитель) регистрирует партнёра через приложение ``Мой налог'' как поставщика услуг.
  \item Требуется подтверждение статуса несовершеннолетнего (паспорт родителя + СНИЛС ребёнка).
  \item Договор с родителями (не с ребёнком---самозанятым) организуется через законного представителя.
\end{enumerate}

\textbf{Важно:} Официальные доходы обязаны быть зарегистрированы в приложении ``Мой налог''. Платежи принимаются только на счёт самозанятого.

\subsection{Договоры с родителями учеников}

Поскольку ученики --- несовершеннолётние, договор заключается с их законными представителями.

\begin{longtable}{p{4cm}p{6cm}}
\caption{Элементы договора с родителями}
\label{tab:contract_elements} \\
\toprule
\textbf{Элемент договора} & \textbf{Содержание} \\
\midrule
\endfirsthead

\toprule
\textbf{Элемент договора} & \textbf{Содержание} \\
\midrule
\endhead

Наименование сторон & 1) [ФИО партнёра] --- самозанятый (представитель: [ФИО родителя]) \\
& 2) [ФИО родителя ученика] --- законный представитель [ФИО ученика] \\
Предмет договора & Оказание образовательных услуг по программам ML, математика, физика \\
Оплата & 750---1000 руб./час/ученик, оплата за месяц вперёд \\
Срок действия & 1 месяц (автоматическое продление) \\
Ответственность & Стороны несут ответственность за нарушение сроков и качества \\
Принятие участия & Родитель подтверждает согласие ребёнка участвовать \\
\bottomrule
\end{longtable}

\textbf{Рекомендуется:} Использовать образец договора, сертифицированный Министерством просвещения РФ.

\subsection{Совместная работа двух партнёров}

Поскольку самозанятость --- индивидуальный статус, каждый партнёр:

\begin{enumerate}
  \item Регистрируется самостоятельно (через родителей).
  \item Получает доход на свой счёт.
  \item Платит налог самостоятельно.
  \item Не может передавать работу третьим лицам.
\end{enumerate}

\textbf{Кооперация:}
\begin{itemize}
  \item \textbf{Разделение аудиторий:} Партнёр 1 --- ML и математика; Партнёр 2 --- физика и математика.
  \item \textbf{Совместные проекты:} Оба участвуют в разработке курсов и маркетинге.
  \item \textbf{Финансы:} Кооперативные расходы (маркетинг, материалы) оплачиваются по очереди или поровну.
\end{itemize}

\subsection{Избегание серых схем}

\begin{enumerate}
  \item \textbf{Только через ``Мой налог'':} Все поступления регистрируются в приложении.
  \item \textbf{Безналичные расчёты:} Только карта/СБП для самозанятых (получение средств автоматически регистрирует доход).
  \item \textbf{Документальное оформление:} Договор с родителями, акт выполненных работ.
  \item \textbf{Отчётность:} Ежемесячная проверка в приложении ``Мой налог''.
\end{enumerate}

\subsection{Перспектива выхода на ИП/ООО}

До 18 лет невозможно зарегистрировать ИП. План перехода:

\begin{table}[H]
\centering
\caption{Этапы правового развития}
\label{tab:legal_timeline}
\begin{tabular}{ll}
\toprule
Возраст & Статус \\
\midrule
16---17 лет & Самозанятые (НПД) \\
18 лет & Индивидуальный предприниматель (упрощёнка 6\%/15\%) \\
${>}$ 20 лет & ООО (в случае привлечения инвестиций или партнёров) \\
\bottomrule
\end{tabular}
\end{table}

% ========================================================================
% МАРКЕТИНГ
% ========================================================================
\section{Маркетинг}

\subsection{Стратегия привлечения учеников}

\subsubsection{Этап 1: Собственные сети (месяцы 1---3)}

\begin{itemize}
  \item \textbf{Семья и друзья:} Прямые рекомендации (80\% успеха в первой очереди).
  \item \textbf{Школьные сети:} Разрешение на раздачу листовок в лицее (бесплатно).
  \item \textbf{Однокурсники:} Рассказ однокурсникам, их родителям.
\end{itemize}

Ожидаемый результат: 5---8 учеников на старте.

\subsubsection{Этап 2: Соцсети (месяцы 4---8)}

\begin{itemize}
  \item \textbf{Telegram-канал:} Еженедельные посты о ML, решение олимпиадных задач, отзывы учеников.
  \item \textbf{Instagram:} Карусели с объяснениями тем, отзывы, живые занятия.
  \item \textbf{YouTube Shorts / VK Video:} 1---2 минутные видео с объяснением.
\end{itemize}

Бюджет: 1\,000---2\,000 руб./мес на таргет (VK, Instagram).

\subsubsection{Этап 3: Реферальная программа (месяцы 5---12)}

\begin{itemize}
  \item \textbf{Бонус 500 руб.} на счёт ученика за приглашение друга.
  \item \textbf{Бонус 500 руб.} на счёт друга (первое занятие бесплатно).
  \item \textbf{Кешбэк 10\%} от рекомендационного дохода на следующий месяц.
\end{itemize}

\subsubsection{Этап 4: Партнёрство с школами (месяцы 9---12)}

\begin{itemize}
  \item Предложить бесплатные мастер-классы по ML в школах.
  \item Стать рекомендованным репетитором в школе.
  \item Участвовать в школьных ярмарках/олимпиадах.
\end{itemize}

\subsection{Брендинг}

\begin{enumerate}
  \item \textbf{Слоган:} ``ML для школьников, не для инженеров''.
  \item \textbf{УТП:} Доступное ML-образование от сверстника, который сам прошёл этот путь.
  \item \textbf{Визуал:} Современный, яркий, понятный даже 10-летнему ребёнку.
\end{enumerate}

\subsection{Минимальный рекламный бюджет}

\begin{table}[H]
\centering
\caption{Маркетинговый бюджет (сумма за месяц)}
\begin{tabular}{lr}
\toprule
Канал & Бюджет, руб. \\
\midrule
Соцсети (таргет) & 1\,500 \\
Печатные листовки & 300 \\
Пробные занятия (бонусы) & 2\,000 \\
YouTube/TG (контент) & 0 (собственное время) \\
\midrule
\textbf{Итого} & \textbf{3\,800} \\
\bottomrule
\end{tabular}
\end{table}

% ========================================================================
% РИСКИ И МИТИГАЦИИ
% ========================================================================
\section{Риски и меры их снижения}

\begin{longtable}{p{4cm}p{4cm}p{6cm}}
\caption{Риски и меры их снижения}
\label{tab:risks} \\
\toprule
\textbf{Риск} & \textbf{Тип} & \textbf{Митигация} \\
\midrule
\endfirsthead

\toprule
\textbf{Риск} & \textbf{Тип} & \textbf{Митигация} \\
\midrule
\endhead

Низкий спрос на старте & Операционный & Пробные занятия бесплатно, реферальные бонусы, работа через школы \\
Перегрузка (выгорание) & Операционный & Чередование групп между партнёрами, план отдыха, 1.5 часа/группа \\
Потеря учеников & Коммерческий & Скидка на первый месяц, гибкая система скидок, подарочные сертификаты \\
Лимит дохода 2,4 млн & Финансовый & Деление дохода между партнёрами, реинвестирование в развитие \\
Сложности с регистрацией & Юридический & Подача документов за 2 недели, помощь родителей \\
Конкуренция со стороны МШП & Стратегический & Доступная цена, персональный подход, ML-специализация \\
Сезонные колебания & Операционный & Летние интенсивы, подготовка к экзаменам, резервный фонд \\
Несвоевременные платежи & Финансовый & Предоплата за месяц, штраф 500 руб. за просрочку \\
Ограничения несовершеннолетних & Юридический & Работа только с родителями, договор через представителей \\
Трудности с курендёрой ML & Продуктовый & Использование проверенных open-source, постепенное усложнение \\
\bottomrule
\end{longtable}

\subsection{Финансовый резерв}

Рекомендуется откладить 20\% от чистой прибыли на резервный фонд. На старте (10\,000 руб.) резерв покрывает 2 месяца эксплуатации.

% ========================================================================
% ФИНАНСОВЫЙ ПРОГНОЗ
% ========================================================================
\section{Финансовый прогноз}

\subsection{Прогноз доходов (год 1)}

\begin{table}[H]
\centering
\caption{Месячный прогноз доходов и прибыли (Год 1)}
\label{tab:year1_forecast}
\begin{tabular}{lrrrrr}
\toprule
Месяц & Доход, руб. & Расходы, руб. & Доходимость & Налог, руб. & Прибыль/партнёр, руб. \\
\midrule
Октябрь & 0 & 8\,000 & 0\% & 0 & $-$4\,000 \\
Ноябрь & 70\,000 & 5\,500 & 92\% & 0 & 32\,250 \\
Декабрь & 125\,000 & 6\,000 & 95\% & 0 & 59\,500 \\
Январь & 140\,000 & 6\,000 & 96\% & 0 & 67\,000 \\
Февраль & 185\,000 & 5\,000 & 97\% & 0 & 90\,000 \\
Март & 250\,000 & 6\,000 & 98\% & 0 & 122\,000 \\
Апрель & 220\,000 & 5\,000 & 98\% & 0 & 107\,500 \\
Май & 190\,000 & 4\,000 & 98\% & 0 & 93\,000 \\
Июнь & 130\,000 & 4\,000 & 97\% & 0 & 63\,000 \\
Июль & 196\,000 & 5\,000 & 97\% & 0 & 95\,500 \\
Август & 45\,000 & 3\,000 & 93\% & 0 & 21\,000 \\
Сентябрь & 215\,000 & 6\,000 & 97\% & 0 & 104\,500 \\
\midrule
\textbf{Итого} & \textbf{1\,766\,000} & \textbf{54\,500} & \textbf{97\%} & \textbf{0} & \textbf{831\,750} \\
\bottomrule
\end{tabular}
\end{table}

\textbf{Пояснение:} Налог в таблице --- 0, так как доход одного партнёра (883\,000 руб.) не превышает порог 400\,000 руб./год для обложения НПД. Чистая прибыль на партнёра после расходов: 831\,750 руб.

\subsection{Прогноз доходов (год 2---3)}

\begin{table}[H]
\centering
\caption{Прогноз доходов --- Год 2 и Год 3}
\label{tab:year23_forecast}
\begin{tabular}{lrrr}
\toprule
Показатель & Год 2 & Год 3 & Рост \\
\midrule
Количество групп (пик) & 12 & 18 \\
Среднее учеников/группа & 6 & 6 \\
Средняя цена, руб./час & 950 & 1\,100 \\
Месячный доход (пик), руб. & 450\,000 & 750\,000 \\
Годовой доход (брутто), руб. & 3\,500\,000 & 5\,500\,000 \\
Доход на партнёра, руб. & 1\,750\,000 & 2\,750\,000 \\
Налог на партнёра, руб. & 54\,000 & 94\,000 \\
Чистый доход на партнёра, руб. & 1\,696\,000 & 2\,656\,000 \\
\bottomrule
\end{tabular}
\end{table}

\textbf{Допущения:}
\begin{enumerate}
  \item Рост цены на 5\% в год (инфляция).
  \item Увеличение количества групп за счёт найма ассистентов (ГПХ) или выхода на ИП статус.
  \item В Год 3 планируется переход на ИП (при 18 лет) --- упрощённая система налогообложения 6\% до 1,5 млн, 15\% свыше.
\end{enumerate}

\subsection{Модель масштабирования до МШП (Год 3---5)}

\begin{table}[H]
\centering
\caption{Модель масштабирования до образовательной компании}
\label{tab:scaling_model}
\begin{tabular}{lrrr}
\toprule
Показатель & Год 3 & Год 4 & Год 5 \\
\midrule
Количество партнёров---основателей & 2 & 5 & 10 \\
Количество групп & 18 & 40 & 100 \\
Среднее учеников/группа & 6 & 6 & 6 \\
Месячный доход (брутто), руб. & 750\,000 & 1\,650\,000 & 3\,500\,000 \\
Годовой доход (брутто), руб. & 5\,500\,000 & 12\,000\,000 & 25\,000\,000 \\
Чистая прибыль (после налогов), руб. & 3\,500\,000 & 7\,500\,000 & 15\,000\,000 \\
\bottomrule
\end{tabular}
\end{table}

\textbf{Ключевые этапы масштабирования:}
\begin{enumerate}
  \item \textbf{Год 2:} Найм ассистентов (старшие школьники, студенты) по ГПХ / договору об образовании.
  \item \textbf{Год 3:} Переход на ИП, создание фирмы-некоммерческой организации для привлечения грантов.
  \item \textbf{Год 4---5:} Создание ООО, привлечение партнёров (преподаватели с долей), развитие франшизы.
\end{enumerate}

% ========================================================================
% ВЫВОДЫ
% ========================================================================
\section*{Выводы}
\addcontentsline{toc}{section}{Выводы}

\begin{enumerate}
  \item \textbf{Реалистичный старт:} Год 1 --- 830\,000 руб./партнёр чистыми. Это реальная и достижимая цель для подростков с опытом и бесплатной площадкой.
  \item \textbf{ML --- ключевая ниша:} Рынок ML-образования для школьников недооценён. Это позволяет выделиться среди конкурентов.
  \item \textbf{Юридические ограничения:} Самозанятость для несовершеннолетних работает, но требует терпения с родителями. Переход на ИП в 18 лет --- критически важный этап.
  \item \textbf{Рост до МШП:} Достижимо за 3---5 лет при правильном подходе к найму и брендингу. Главное --- выжить в Год 1 и наладить стабильный доход.
  \item \textbf{Риски:} Главное --- выгорание и потеря мотивации. Рекомендуется чётко планировать нагрузку и использовать автоматизацию.
\end{enumerate}

\subsection*{Рекомендации}

\begin{enumerate}
  \item Начать с 2---3 групп, а не сразу с максимальной загрузкой.
  \item Инвестировать первые 3 месяца в контент и обратную связь, а не в рекламу.
  \item Регулярно проверять границы дохода по НПД (2,4 млн руб./год на партнёра).
  \item Строить систему рекомендаций: каждый новый ученик --- это 3 потенциальных клиента.
  \item Не экономить на качестве: репутация подросткового репетитора строится быстро, но теряется быстрее.
\end{enumerate}

% ========================================================================
% ПРИЛОЖЕНИЯ
% ========================================================================

\newpage
\section*{Приложение А: Образец договора с родителями}
\addcontentsline{toc}{section}{Приложение А: Образец договора}

\textbf{ДОГОВОР № \_\_\_\_\_\_ от \_\_\_\_\_\_\_\_\_\_\_\_\_\_\_\_}

{\ragged_left Между: \par}
{\raggedright
\textbf{[ФИО партнёра]}, действующий(ая) от имени которого --- [ФИО родителя] (законный представитель), с одной стороны, \\
\textbf{[ФИО родителя ученика]}, с другой стороны, \\
(далее совместно --- Стороны), \\
}

\textbf{1. Предмет договора:} Оказание образовательных услуг по программам: машинное обучение, математика, физика для учащегося [ФИО] [класс] класса.

\textbf{2. Порядок оплаты:} 750---1000 руб./час/ученик. Оплата --- 1 раз в месяц (до 5 числа) банковской картой самозанятого через СБП.

\textbf{3. График занятий:} [дни недели], [время], [формат: онлайн/офлайн].

\textbf{4. Ответственность:} Стороны несут ответственность в соответствии с ФЗ.

\textbf{5. Прочие условия:} Ученик обязуется выполнять домашние задания. Партнёр---обязуется предоставлять качественные материалы.

{\raggedleft Подписи сторон: \par}
\_\_\_\_\_\_\_\_\_\_\_\_/ \_\_\_\_\_\_\_\_\_\_\_\_ \hrulefill

(родитель---участник / репетитор)

\newpage
\section*{Приложение Б: Чек-лист старта}
\addcontentsline{toc}{section}{Приложение Б: Чек-лист старта}

\begin{longtable}{p{2cm}p{8cm}p{2cm}}
\caption{Чек-лист старта проекта}
\label{tab:checklist} \\
\toprule
№ & Действие & Статус \\
\midrule
\endfirsthead

\toprule
№ & Действие & Статус \\
\midrule
\endhead

1 & Регистрация в ``Мой налог'' (через родителей) & \_\_\,/\\_ \\
2 & Разработка программы ML (10 уроков) & \_\_\,/\\_ \\
3 & Разработка программы математики & \_\_\,/\\_ \\
4 & Создание видео/превью для рекламы & \_\_\,/\\_ \\
5 & Раздача листовок в лицее & \_\_\,/\\_ \\
6 & Запуск реферальной программы & \_\_\,/\\_ \\
7 & Первые 3 ученика & \_\_\,/\\_ \\
8 & Первая выплата через СБП & \_\_\,/\\_ \\
9 & Отзыв первого ученика & \_\_\,/\\_ \\
10 & Запуск второй группы & \_\_\,/\\_ \\
\bottomrule
\end{longtable}

\end{document}